\documentclass[a4paper]{article}
\usepackage[french]{babel}
\usepackage[utf8]{inputenc}
\usepackage[T1]{fontenc}
\usepackage{hyperref}
\usepackage{graphicx}
\usepackage{appendix}
\usepackage{caption}
\usepackage{pdfpages}
\usepackage{mathtools}
\usepackage{cite}
\usepackage{array}
\usepackage{multirow}
\usepackage{listings}
\lstset{language=Python}
\begin{document}

\begin{titlepage}
\newcommand{\HRule}{\rule{\linewidth}{0.5mm}} 
\center 
 
%----------------------------------------------------------------------------------------
%	HEADING SECTIONS
%----------------------------------------------------------------------------------------

\textsc{\LARGE Université de Bordeaux}\\[1.5cm] 
\textsc{\Large Projet Métaheuristique}\\[0.5cm] 
\textsc{\large M1 Informatique}\\[0.5cm] 

%----------------------------------------------------------------------------------------
%	TITLE SECTION
%----------------------------------------------------------------------------------------

\HRule \\[0.4cm]
{ \huge \bfseries Partitionnement de graphe en classes}\\[0.4cm] % Title of your document
\HRule \\[1.5cm]
 
%----------------------------------------------------------------------------------------
%	AUTHOR SECTION
%----------------------------------------------------------------------------------------

\begin{minipage}{0.4\textwidth}
\begin{flushleft} \large
\emph{Auteurs : }\\
Norbert \textsc{Feron},\\
Baptiste \textsc{Masset}
\end{flushleft}
\end{minipage}
~
\begin{minipage}{0.4\textwidth}
\begin{flushright} \large
\emph{Encadrant : } \\
Marc Michel \textsc{Corsini} 
\end{flushright}
\end{minipage}\\[2cm]

%----------------------------------------------------------------------------------------
%	DATE SECTION
%----------------------------------------------------------------------------------------

{\large \today}\\[2cm]

%----------------------------------------------------------------------------------------
%	LOGO SECTION
%----------------------------------------------------------------------------------------

\includegraphics{img/logo.png}\\[1cm]
 
%----------------------------------------------------------------------------------------

\vfill
\end{titlepage}
%----------------------------------------------------------------------------------------

\tableofcontents

\newpage
\section{Introduction}
Dans le cadre du cours de recherche opérationnel de master 1 informatique a l'université de Bordeaux, ce projet a pour but de mettre en pratique des algorithmes et méta-algorithme, tel que la descente de gradient, le Recuit simulé ou encore la méthode Tabou. Pour cela nous disposons de graphes pondérés de différentes taille (de 5-7 sommets a plus de milles), sur lesquelles nous allons tester et comparé les résultats de nos différents algorithme appliqué au problème de partitionnement de graphe en classes à peu près équitables.

\subsection{Le problème}
Notre problème: partitionner un graphe en K classes en minimisant le poids interclasses. Cela consistes à placer chaque sommet du graphe dans une des classe de manières à réduire le somme du poids des arrêtes n'appartenant pas à la même classe.

\subsection{Nombre de classes}
Nous avons choisi de laisser le nombre de classes libre, l'utilisateur lance donc les algorithmes en donnant comme entrée un graphe pondéré et un nombre de classe. 

\subsection{Définition de l'équitable}

Nous avons choisi de ne pas implémenter la notion  d'équilibre entre les classes dans nos méta-algorithmes, l'utilisateur dispose d'un variable supplémentaire appelé delta\_max; elle correspond à l'écart maximum entre les classes pour qu'une solution soit considérée comme valide et puisse être retournée en sortie des différents algorithmes utilisés dans ce projet.


\section{Méthode d’énumération}
La solution la plus simple consiste à énumérer toutes les solutions valides et de calculer la somme du poids des arrêtes interclasses, les solutions avec la somme des poids plus faible sont retournées en résultats de l'algorithme d'énumération. Les solutions retournées sont donc toujours les solutions optimales, cependant la puissance de nos calculateurs fait défaut et en dépassant un nombre de sommets ou en demandant un partitionnement avec plusieurs classes, nos ordinateurs ne dispose pas d'un puissance de calcul suffisante et l'énumération des solutions devient impossible.

\subsection{Test et limite}

\section{Descente de Gradient}

\section{Recuit Simulé}

\section{Méthode Tabou}

\section{Comparaison des résultats}

\section{Conclusion}

\end{document}